\chapter{El proyecto de un vistazo}

Este libro documenta, paso a paso, cómo se construyó la simulación de monitoreo de incendios:
desde la instalación del software hasta la misión autónoma en la que un cuadrirrotor
Tarot Ironman 650 recorre un bosque, encuentra un incendio con una cámara y una red neuronal,
envía sus coordenadas GPS y lo orbita. Todo el desarrollo se hizo entre el 21 y el 23 de
octubre de 2026 en el espacio de trabajo \cod{~/ros2_ws}.

\section{Qué hace la simulación}

\begin{itemize}
  \item \textbf{Planta.} Gazebo Harmonic simula la física del dron: cuerpo rígido, cuatro
        motores de primer orden con empuje $f_i = k_f\omega_i^2$, una cámara, una IMU y un
        LiDAR 3D. Hay dos vehículos: un cuadrirrotor genérico en ``+'' y el Tarot 650 en ``X'',
        cuya geometría viene de su CAD.
  \item \textbf{Control.} Tres alternativas intercambiables: un PD en cascada (línea base),
        el super-twisting adaptativo (STA) con tres discretizaciones (explícita,
        semi-implícita e implícita) y el piloto automático ArduCopter 4.7.2 en SITL.
  \item \textbf{Mundos.} Tres escenarios: \cod{incendio} (plano), \cod{bosque} (85 árboles) y
        \cod{bosque_grande}: un valle de $420\times420$~m rodeado de montañas, con 1700 árboles,
        140 rocas, 21 animales, cielo y un incendio con humo.
  \item \textbf{Percepción.} YOLO26n ajustado con 6000 imágenes del simulador detecta
        \emph{fuego} y \emph{humo} en la cámara del dron.
  \item \textbf{Autonomía.} El nodo \cod{busqueda_incendio} despega, barre el bosque en
        serpentina siguiendo el relieve, evita obstáculos con el LiDAR, estima la posición del
        foco, lo confirma con las llamas, avisa con las coordenadas GPS y lo orbita.
\end{itemize}

\section{Arquitectura}

La figura~\ref{fig:arquitectura} muestra los nodos y los tópicos principales. La cadena completa es
\emph{imagen $\rightarrow$ detección $\rightarrow$ estimación $\rightarrow$
referencia $\rightarrow$ control}. El planificador nunca usa la posición verdadera del foco:
solo lo que reporta el detector.

\begin{figure}[H]
\centering
\begin{tikzpicture}[
  nodo/.style={draw,rounded corners,fill=#1,minimum width=3.1cm,minimum height=1cm,align=center,font=\small},
  nodo/.default=white,
  flecha/.style={-{Stealth[length=2.5mm]},thick},
  et/.style={font=\scriptsize\ttfamily,fill=white,inner sep=1pt}]
  \node[nodo=orange!15,minimum width=4.2cm,minimum height=2.2cm] (gz) at (0,0)
       {\textbf{Gazebo Harmonic}\\ física, motores, cámara,\\ IMU, LiDAR, mundo};
  \node[nodo=black!5] (puente) at (0,-2.4) {\cod{ros_gz_bridge}\\(\cod{puente.yaml})};
  \node[nodo=blue!10] (yolo) at (-5,-4.6) {\cod{detector_yolo}\\YOLO26n (GPU)};
  \node[nodo=green!12,minimum width=4cm] (bus) at (0,-6.6) {\cod{busqueda_incendio}\\planificador autónomo};
  \node[nodo=red!10] (ctl) at (5,-4.6) {\cod{controlador_sta}\\ o \cod{controlador_pd}};
  \node[nodo=red!10] (ap) at (5,-8.6) {\cod{mision_ardupilot}\\+ ArduCopter SITL};
  \node[nodo=black!5] (fauna) at (-6.2,0) {\cod{fauna}\\(gz-transport)};
  \node[nodo=black!5] (viento) at (6.2,0) {\cod{viento}};
  \node[nodo=yellow!15] (alerta) at (-4.5,-8.6) {aviso GPS\\\cod{/cuadrirotor/incendio/*}};

  \draw[flecha] (gz) -- (puente);
  \draw[flecha] (puente) -| node[et,pos=0.75]{camara} (yolo);
  \draw[flecha] (yolo) |- node[et,pos=0.3]{detecciones} (bus);
  \draw[flecha] (puente) -- node[et]{odometria, lidar/points} (bus);
  \draw[flecha] (bus) -| node[et,pos=0.3]{consigna} (ctl);
  \draw[flecha] (bus) -| node[et,pos=0.3]{consigna} (ap);
  \draw[flecha] (ctl) |- node[et,pos=0.75]{motores} (puente);
  \draw[flecha] (ap.east) -- ++(2.6,0) |- node[et,pos=0.25,rotate=90]{MAVLink / JSON} ([yshift=-7mm]gz.east);
  \draw[flecha] (fauna) -- node[et,above=1pt]{set\_pose\_vector} (gz);
  \draw[flecha] ([yshift=3mm]viento.west) -- node[et,above=1pt]{viento} ([yshift=3mm]gz.east);
  \draw[flecha] (bus.south) |- (alerta.east);
\end{tikzpicture}
\caption{Nodos y tópicos de la misión de búsqueda. Con \cod{controlador:=sta} o \cod{pd} la
consigna la sigue el controlador de ROS~2; con \cod{controlador:=ardupilot}, la reenvía
\cod{mision_ardupilot} a ArduCopter por MAVLink.}
\label{fig:arquitectura}
\end{figure}

\section{Convenciones}

\begin{itemize}
  \item \textbf{Marco del mundo:} ENU (\emph{x} al este, \emph{y} al norte, \emph{z} arriba), en
        metros. ArduPilot usa NED: $N = y$, $E = x$, $D = -z$ y la guiñada NED es
        $\pi/2 - \psi$.
  \item \textbf{Marco del cuerpo:} \emph{x} al frente, \emph{y} a la izquierda, \emph{z} arriba.
  \item El modelo del dron se llama siempre \cod{cuadrirotor} dentro de Gazebo, sin importar
        el vehículo; así los tópicos no cambian.
  \item Rutas: el paquete principal está en \cod{~/ros2_ws/src/cuadrirotor_sim}; salvo que se
        diga otra cosa, las rutas relativas son respecto a él.
\end{itemize}

\section{Equipo usado}

Ubuntu 24.04 (Noble), ROS~2 Jazzy, Gazebo Harmonic (gz-sim 8), ArduCopter 4.7.2 y una GPU
NVIDIA RTX 5060 Ti (YOLO corre en la GPU y Gazebo renderiza con EGL de NVIDIA). Con este equipo
el mundo grande corre a tiempo real (factor $\approx 1.0$; $\approx 0.8$ con el LiDAR).
